\section{Experimental Setup}
\label{sec:setup}

\subsection{Implementation Environment and Compute Infrastructure}
All experiments, model training pipelines, and data audit procedures were implemented in Python 3.11 using the PyTorch 2.4 deep learning framework. Hardware training was conducted on high-performance compute nodes equipped with an NVIDIA RTX A6000 GPU (48\,GB GDDR6 VRAM, CUDA 12.2) and an AMD EPYC 7763 64-core processor running Ubuntu 22.04 LTS.

For edge deployment profiling, models were exported to the Open Neural Network Exchange (ONNX) format and benchmarked using ONNX Runtime 1.19 across three representative computing platforms:
\begin{enumerate}
    \item \textbf{High-Performance Workstation CPU:} Intel Core i7-12700K (12 physical cores, 20 threads, base frequency 3.60\,GHz, 32\,GB DDR5 memory).
    \item \textbf{Low-Cost Agricultural Edge SBC:} Raspberry Pi 4 Model B (Broadcom BCM2711 SoC, Quad-core ARM Cortex-A72 64-bit @ 1.5\,GHz, 4\,GB LPDDR4-3200 SDRAM, running Raspberry Pi OS 64-bit).
    \item \textbf{Embedded Neural Accelerator:} NVIDIA Jetson Nano Developer Kit (128-core NVIDIA Maxwell GPU, Quad-core ARM Cortex-A57 @ 1.43\,GHz, 4\,GB 64-bit LPDDR4, running JetPack 4.6.1 / TensorRT).
\end{enumerate}

\subsection{Hyperparameter Specifications and Optimization Budget}
To ensure rigorous and equitable benchmarking, all baseline architectures and LitchiHybridNet variants were trained under strictly identical optimization budgets and data splits. Table~\ref{tab:hyperparams} enumerates the complete hyperparameter configuration.

\begin{table}[t]
\centering
\caption{Hyperparameter Specifications for Model Training and Gabor Layer}
\label{tab:hyperparams}
\begin{tabular}{ll}
\toprule
\textbf{Hyperparameter} & \textbf{Value / Specification} \\
\midrule
Base Architecture & MobileNetV3-Large (ImageNet-1k pre-trained) \\
Input Image Dimensions & $224 \times 224 \times 3$ RGB \\
Optimization Algorithm & AdamW ($\beta_1 = 0.9, \beta_2 = 0.999$) \\
Initial Learning Rate & $1.0 \times 10^{-3}$ \\
Learning Rate Schedule & Cosine Annealing ($\eta_{\min} = 1.0 \times 10^{-6}$) \\
Weight Decay Penalty & $1.0 \times 10^{-2}$ \\
Batch Size & 32 \\
Training Duration & 50 Epochs (Early stopping patience = 10) \\
Label Smoothing ($\epsilon$) & 0.1 \\
Dropout Probability ($p$) & 0.3 \\
Gabor Bank Channel Count ($K$) & 24 \\
Gabor Spatial Kernel Size & $7 \times 7$ \\
Initial Gabor Orientations ($\theta$) & 8 Angles: $\{0, \frac{\pi}{8}, \frac{\pi}{4}, \frac{3\pi}{8}, \frac{\pi}{2}, \frac{5\pi}{8}, \frac{3\pi}{4}, \frac{7\pi}{8}\}$ \\
Initial Radial Frequencies ($F$) & 3 Frequencies: $\{0.05, 0.15, 0.25\}$ cycles/pixel \\
Initial Envelope Aspect Ratio ($\gamma$) & 1.0 (Isotropic envelope) \\
Quantization Paradigm & Post-Training Dynamic INT8 via ONNX Runtime \\
\bottomrule
\end{tabular}
\end{table}


All networks were optimized using the AdamW algorithm ($\beta_1 = 0.9, \beta_2 = 0.999$, weight decay $\lambda = 1.0 \times 10^{-2}$) with an initial learning rate $\eta_0 = 1.0 \times 10^{-3}$. The learning rate was modulated using a Cosine Annealing schedule decaying smoothly to $\eta_{\min} = 1.0 \times 10^{-6}$ over $E = 50$ epochs. The batch size was fixed at 32. Early stopping based on validation Macro-F1 was enforced with a patience of 10 epochs.

\subsection{Comparative Baseline Architectures}
We evaluate LitchiHybridNet against a diverse suite of 8 competitive baselines spanning classical machine learning, lightweight convolutional models, modern high-capacity backbones, vision transformers, and graph neural networks:
\begin{enumerate}
    \item \textbf{Classical Gabor + SVM:} A traditional hand-crafted baseline combining multi-scale Gabor filter bank texture energy features with a Radial Basis Function (RBF) Support Vector Machine ($C=10, \gamma=\text{scale}$).
    \item \textbf{ShuffleNetV2 1.0$\times$ \cite{ma2018shufflenet}:} An ultra-lightweight architecture (2.28\,M parameters) designed around channel split and channel shuffle operations.
    \item \textbf{MobileNetV2 \cite{sandler2018mobilenetv2}:} An inverted residual CNN with linear bottlenecks (3.50\,M parameters).
    \item \textbf{MobileNetV3-Large \cite{howard2019searching}:} The direct backbone baseline (5.48\,M parameters) incorporating hardware-aware NAS and Squeeze-and-Excitation attention.
    \item \textbf{ResNet-50 \cite{he2016deep}:} A canonical 50-layer residual network with bottleneck blocks (25.56\,M parameters).
    \item \textbf{EfficientNet-B0 \cite{tan2019efficientnet}:} A compound-scaled CNN optimizing depth, width, and resolution (5.29\,M parameters).
    \item \textbf{Swin-Transformer-Tiny (Swin-T) \cite{liu2021swin}:} A hierarchical Vision Transformer utilizing shifted local windows (28.29\,M parameters).
    \item \textbf{ConvNeXt-Tiny \cite{liu2022convnet}:} A modernized pure convolutional network designed to match transformer scaling (28.60\,M parameters).
    \item \textbf{LitchiChebNet \cite{litchi_dataset_2022}:} The reported spectral graph convolutional benchmark utilizing Chebyshev polynomials over segmented superpixel boundaries.
\end{enumerate}

\subsection{Performance Metrics and Statistical Protocol}
Evaluation relies on standard diagnostic metrics computed over the held-out test split of \testEvaluatedSamples{} images:
\begin{itemize}
    \item \textbf{Top-1 Classification Accuracy:} Proportion of correctly predicted samples across all categories.
    \item \textbf{Macro-Averaged Precision, Recall, and F1-Score:} Arithmetic means of per-class metrics, ensuring that minority and majority classes are weighted equally:
    \begin{equation}
        \text{Macro-F1} = \frac{1}{C} \sum_{c=1}^C \frac{2 \cdot P_c \cdot R_c}{P_c + R_c}
    \end{equation}
    \item \textbf{Weighted-Averaged F1-Score:} Per-class F1-scores weighted by class sample support.
    \item \textbf{Matthews Correlation Coefficient (MCC):} A balanced measure robust to class size disparities:
    \begin{equation}
        \text{MCC} = \frac{c \cdot s - \sum_k p_k t_k}{\sqrt{\left( s^2 - \sum_k p_k^2 \right) \left( s^2 - \sum_k t_k^2 \right)}}
    \end{equation}
    where $c$ is the total correct predictions, $s$ is the total samples, and $p_k, t_k$ denote predicted and true marginal totals.
\end{itemize}

To assess statistical significance, all experiments were replicated across 5 independent random seeds $\{42, 123, 456, 789, 999\}$. We report mean values accompanied by standard deviations ($\pm \sigma$). Pairwise model differences are subjected to McNemar's test for classification discordance on the 1,665 test samples, and the non-parametric Wilcoxon signed-rank test across cross-validation runs.

\subsection{Environmental Corruption Robustness Protocol}
To evaluate diagnostic reliability under adverse orchard conditions, we simulate five distinct field corruptions adapted from the Hendrycks robustness protocol \cite{hendrycks2019benchmarking} across five severity levels ($s \in \{1, 2, 3, 4, 5\}$):
\begin{enumerate}
    \item \textbf{Optical Motion Blur:} Simulates camera shake and foliage flutter via linear motion kernels of length $L \in \{5, 9, 15, 21, 27\}$ pixels.
    \item \textbf{Gaussian Sensor Noise:} Simulates low-light sensor noise with variance $\sigma_n^2 \in \{0.02, 0.05, 0.08, 0.12, 0.18\}$.
    \item \textbf{Monsoon Rain Simulation:} Synthetic rain streaks and lens droplet splatters generated by semi-transparent linear blurs and refractive circular discs.
    \item \textbf{Solar Glare / Overexposure:} Non-linear illumination flares and localized intensity saturation ($I_{\text{new}} = I^\gamma$, $\gamma \in [0.85, 0.45]$).
    \item \textbf{Atmospheric Fog / Contrast Loss:} Haze and scattering simulated via the Koschmieder optical transmission model with scattering coefficients $\beta \in [0.5, 2.5]$.
\end{enumerate}
